\begin{figure}[htbp]
\centering
\begin{tikzpicture}[
  font=\small,
  main/.style={draw=black!65,rounded corners=2pt,fill=black!3,
    align=center,text width=6.2cm,minimum height=1.12cm,inner sep=5pt},
  scalar/.style={draw=blue!45!black,rounded corners=2pt,fill=blue!4,
    align=center,text width=5.2cm,minimum height=1.3cm,inner sep=5pt},
  use/.style={-{Stealth[length=2.2mm]},draw=black!70,line width=.65pt}
]
\node[main] (cone) at (0,0)
  {Cone formulation\\Section~\ref{sec:01-cones}};
\node[main] (projection) at (0,-1.65)
  {Normalized Gaussian projections\\and entropy estimate\\
   Sections~\ref{sec:02-projections}--\ref{sec:03-entropy}};
\node[main] (quadratic) at (0,-3.3)
  {Quadratic identities and\\layer error bounds\\
   Sections~\ref{sec:04-quadratic}--\ref{sec:05-layers}};
\node[main] (linear) at (0,-4.95)
  {Linear-field estimate and\\equality characterization\\
   Section~\ref{sec:06-linear-equality}};
\node[align=center,font=\small\bfseries] at (9,.35)
  {Independent scalar input};
\node[scalar] (profiles) at (9,-1)
  {Scalar profile bounds\\
   Stated in Section~\ref{sec:01-cones};\\
   proved in Appendix~\ref{sec:07-scalar}};
\node[scalar] (segments) at (9,-4.95)
  {Segment variance estimate\\Appendix~\ref{sec:08-segments}};
\draw[use] (cone) -- (projection);
\draw[use] (projection) -- (quadratic);
\draw[use] (quadratic) -- (linear);
\draw[use] (profiles) -- (segments);
\draw[use] (profiles.west) -- (projection.east);
\draw[use] (profiles.south west) -| (5.2,-3.3) -- (quadratic.east);
\draw[use] (segments.west) -- (linear.east);
\end{tikzpicture}
\caption{Principal logical dependencies. Arrows indicate where inputs are
used, rather than the order of presentation. The scalar branch is independent
of the cone and dimension; its proofs follow the geometric argument.}
\label{fig:proof-map}
\end{figure}
